\documentclass[10pt]{article}

\usepackage[utf8]{inputenc}

\usepackage[T1]{fontenc}

\usepackage{amsmath}

\usepackage{amsfonts}

\usepackage{amssymb}

\usepackage[version=4]{mhchem}

\usepackage{stmaryrd}

\usepackage{graphicx}

\usepackage[export]{adjustbox}

\graphicspath{ {./images/} }

\usepackage{mathrsfs}



\begin{document}

Первонаqально Э. была определена А. Н. Колмогоро вым несколько иначе (см. [1]); приведенный выше вариант был дан позже (см. [2]). В основном случае aneриодического автомор физма пространства Лебега определения в конечном счете - эквивалентны [3].



Оказывается, что $h\left(S^{n}\right)=n h(S)$, а если $S$ - автоморфизм, то $h\left(S^{-1}\right)=h(S)$. Поэтому Ә. каскада $\left\{S^{n}\right\}$ естественно считать $h(S)$. Для измеримого потока $\left\{S_{t}\right\}$ оказывается, что $h\left(S_{t}\right)=|t| h\left(S_{1}\right)$. Поэтому $Ә$. потока естественно считать $h\left(S_{1}\right)$. Несколько иначе определяется Э. для других групп преобразований с инвариантной мерой (она уже не сводится к Ә. одного из преобразований, входящих в эту группу; см. [5], [6]). Имеется модификация Э. для случая бесконечной инвариантной меры [7]; еце одна модификация - $A$ ә н т р о п и я (где $A=\left\{k_{n}\right\}-$ возрастающая последовательность натуральных чисел) - получается при замене $\xi_{S}^{n}$ на



\[

S-k_{1} \xi \vee \ldots \vee S^{-k_{n}} \xi

\]



и $\lim$ на $\varlimsup$ im (см. [8]).



Э. является инвариантом метрического изоморфизма динамич. систем, принципиально отличным от известных ранее инвариантов, в основном связанных со спектром динамической системы. В частности, с помощью әнтропии Бернулли автоморфизмов (см. [1]) впервые установлено существование неизоморфных әргодич. систем с одинаковым непрерывным спектром (что контрастирует с ситуацией для дискретного спектра). В более широком плане роль Э. связана с тем, что вместе с ней в эргодич. теории возникло новое направление - ә н т р о п и й н а я т е о р и я динамич. систем (см. [3], [4] и эргодическая теория).



Э. дает нек-рую среднюю характеристику скорости перемешивания множества малой меры (точнее, набора таковых, образующих разбиение). Наряду с этой «глобальной» ролью Э. играет и «локальную» роль, к-рая устанавливается э р г о д и ч е с к о й т е о р е м о і̆ Б р ей ма на (индивид у а льная э ргодическая те о рем а теории информ ац и и): для эргодического $S$ при почти всех $x$



\[

\frac{1}{n}\left|\lg \mu\left(C_{\xi_{S}}(x)\right)\right| \rightarrow h(S, \xi) \text { при } n \longrightarrow \infty \text {, }

\]



где $C_{\eta}(x)$ - элемент разбиения $\eta$, содержащий $x$, а логарифм берется по тому же основанию, что и в определении $H$ (см. [9], [4]). (Теорема Бреймана верна для $\xi$ с $H(\xi)<\infty[10]$, но, вообще говоря, неверна для счетных $\xi$ с $H(\xi)=\infty[11]$; имеются также варианты с неәргодическим $S$ (см. [4], [12]) и бесконечной $\mu$ [13]. Более слабое утверждение о сходимости в смысле $L_{1}$ доказано для нек-рого общего класса групп преобразований [6].



Для гладких динамич. систем с гладкой инвариантной мерой установлена связь межнду Ә. и Ляпунова характеристическими показателями уравнений в вариациях (см. [14] - [16]).



Название «Ә.» объясняется аналогией Ә. в теории динамич. систем с Ә. в информации теории и статистич. физике, вплоть до совпадения этих Ә. в нек-рых примерах (см., напр. [4], [17]). Аналогия со статистич. физикой явилась одним из стимулов введения в әргодич. теорию (уже не в чисто метрич. контексте, а для топологических диналических систем) новых понятий - «гиббсовских мер», «топологического давления» (аналог свободной энергии) и «вариационного принципа» для последнего (см. лит. при статьях У-система, Топологическая энтропия).



Jum.: [1] K о л м о г о p o в А. Н., "Докл. АН CCCP", 1958 , т. $119, N_{2} 5$, c. $861-64 ; 1959$, т. 124 , № 4, c. $754-55$ [3] Р о х л и н В. А., "Успехи матем. наук», 1967 , т. 22, в. 5, c. 3-56; [4] Б и л лй "г с ле й П., Эргодическая теория и информация, пер. с англ., М., 1969; [5] С а ф о н о в А. В.,



\includegraphics[max width=\textwidth, center]{2023_07_09_3c8c7070567704f98751g-1(2)}

37; [7] K r e n g e l U., "Z. Wahrscheinlichkeitstheor. und verw. Geb.», 1967, Bd 7, № 3, S. 161-81; [8] К у ш н и p е н к о А. Г., "Успехи матем. наук", 1967, т. 22, B. 5, c. 37-65; [9] B r e i-

m a n L., "Ann. Math. Statistics", 1957, v. 28, № 3, p. $809-11$;



\includegraphics[max width=\textwidth, center]{2023_07_09_3c8c7070567704f98751g-1}

v. 32 , № 3, p. $612-14$; [11] П и к е Ј ь Б. С., «Проблемы передачи информации», 1976 , т. 12, № 2, с. 98-103; [12] I o n e sc u-T u l c e a A., "Arkiv mat.", 1961, Bd 4, H. 2-3, № 18, scheinlicheith $1 \mathrm{~m} \mathrm{k}$ o $\mathrm{L}$. M., S u che ston L.' "Z. Wahr $35 ;$ [14] М и л л и о н щ и к о в В. М., «Дифференциальные уравнения», 1976, т. 12, № 12, с. $2188-92$; [15] П е с и н Я. Б., "Успехи мaтем. наук», 1977, T. 32, в. 4, c. 55-112; [16] M a-

$\tilde{n}$ é R., «Ergodic theory and dynamical systems», 1981, v. 1 , № 1 , p. 95-102; [17] R o b ins on D. W., R u ell'e D., "Ciommun. math.-phys.», 1967, v. 5, № 4, p. 288-300. Д.. В. Аносов.



\&-ЭНТРОПИЯ множества в метрическом пространстве - двоичный логарифм наименьшего числа точек $\varepsilon$-сети для этого множества. Иначе говоря, $\varepsilon-Э$. $\mathscr{H}_{\varepsilon}(C, X)$ множества $C$, лежащего в метрич. пространстве $(X, \rho)$, равна.



\[

\mathscr{H}_{\varepsilon}(C, X)=\log _{2} N_{\varepsilon}(C, X),

\]



где



\[

N_{\varepsilon}(C, X)=\inf \left\{n \mid \exists x_{1}, \ldots, x_{n}, x_{i} \in X:\right.

\]



$\left.C \subset U_{i=1}^{n} B\left(x_{i}, \varepsilon\right)\right\}$, a $B(\zeta, \alpha)=\{x \in X \mid \rho(x, \zeta) \leqslant \alpha\}-$ шар с центром в точке $\zeta$ радиуса $\alpha$. Определение величины $\mathscr{H}_{\varepsilon}(C, X)$ было дано А. Н. Колмогоровым [1], исходившим от нек-рых идей и определений информации теории. Величину $\mathscr{H}_{\varepsilon}(C, X)$ наз. также о т н ос и т е льно й $\varepsilon-э$ н т р опи е й; она зависит от пространства $X$, в к-ром находится множество $C$, т. е. от метрич. расширения множества $C$. Величина $\mathscr{H}_{\varepsilon}(C)=\inf \mathscr{H}_{\varepsilon}(C, X)$,



где нижняя грань берется по всем метрич. расширениям $X$ пространства $C$, наз. а б с о л ю т но й $\varepsilon-\ni$ н тр о пи и е й множества $C$. Величина $\mathscr{H}_{\varepsilon}(C)$ допускает и прямое определение (также предложенное А. Н. Колмогоровым [1]): для метрич. пространства $C$ величина $\mathscr{H}_{\varepsilon}(C)$ есть двоичный логарифм мощности $N_{\varepsilon}(C)$ наиболее әкономного (по числу множеств) $2 \varepsilon$-покрытия $C$.



\includegraphics[max width=\textwidth, center]{2023_07_09_3c8c7070567704f98751g-1(1)}

если $C_{i} \subset C, \bigcup_{i=1}^{n} C_{i}=C$ и диаметр каждого множества не превосходит $2 \varepsilon$. Доказано [2] минимальное свойство абсолютной $\varepsilon-Э$. Величины $N_{\varepsilon}(C), N_{\varepsilon}(C, X)$ суть величины, обратные к поперечникам $\varepsilon^{N}(C)$ и $\varepsilon_{N}(C, X)$, характеризующим наилучшие возможности восстановления элементов из $C$ с помощью таблицы из $N$ әлементов и наилучшие возможности аппроксимации $\mathrm{N}$ точечными множествами. Исследование асимптотики є-Ә. различных функциональных классов составило специальный раздел теории приближений.



Лum.: [1] К о л м о г о p о в А. Н., "Докл. АН СССР», 1956, т. 108, № 3, с. 385-88; [2] В и т у ї к и н А. Г., Оценка сложности задачи табулирования, М., $1959 ;[3]$ К о л м о г оp о в А. Н., Т и х о м и р о в В. М., "Успехи матем. наук",

1959, т. 14. в. 2, с. 3-86.



ЭПИДЕМИИ ПРОЦЕСС - случайный процесс, служащий математич. моделью распространения какойлибо эпидемии. Одна из простейших таких моделей описывается марковским процессом с непрерывным временем, состояния к-рого в момент $t$ описываются числом $\mu_{1}(t)$ больных особей и числом $\mu_{2}(t)$ восприимчивых особей. Если $\mu_{1}(t)=m, \mu_{2}(t)=n$, то за время $(t, t+\Delta t), \Delta t \rightarrow 0$, вероятности переходов определяются следующим образом: $(m, n) \rightarrow(m+1, n-1)$ с вероятностью $\lambda m n \Delta t+o(\Delta t) ; \quad(m, n) \rightarrow(m-1, n)$ с вероятностью $\mu m \Delta t+o(\Delta t)$. В этом случае производящая, функция



\[

F(t ; x, y)=\mathrm{E} x^{\mu_{1}(t)} y^{\mu_{2}(t)}

\]



удовлетворяет диффференциальному уравнению



\[

\frac{\partial F}{\partial t}=\lambda\left(x^{2}-x y\right) \frac{\partial^{2} F}{\partial x \partial y}+\mu(1-x) \frac{\partial F}{\partial x} .

\]



Б. А. Севастъянов.





\end{document}